\documentstyle[% 


righttag% 


,11pt% 


,amssymb% 


,russian%               remove in Eng. 


,izvru%                 Set izven% in Eng. 


]{amsart} 


 


%


\newcommand{\mr}[1]{\mathrm{#1}}


\newcommand{\tts}[1]{} 


 


%\hoffset=-15mm%                  For Eng. 


\setcounter{page}{69} 


\god{2001} 


\nomer{10~(473)} %                        Remove in Eng. 


%\nomer{Vol.\,45, No.\,10}%               Set in Eng. 


%\str{\pageref{begin}--\pageref{end}}%  Set in Eng. 


\udk{514.174.5} 


 


\author{\it В.Г.\,Покровский} 


% NB:   ^^ remove in Eng. 


\title{О разрезании многомерного параллелепипеда\\ на параллелепипеды} 


 


\date{} 


%\pagestyle{myheadings}%         Set in Eng. 


\pagestyle{plain} %              Remove in Eng. 


\begin{document} 


%\thispagestyle{plain}%          Set in Eng. 


\maketitle 


\label{begin} 


 


 


Разрезания прямоугольников на прямоугольники,  


а затем и многомерных параллелепипедов на параллелепипеды изучались  


в ряде работ. В частности, в [1], [2] были получены  


линейные соотношения между ребрами большого и составляющих его  


параллелепипедов: в [1] --- для параллелепипедов  


с целочисленными ребрами, в [2] --- для параллелепипедов  


с произвольными ребрами.  


 


Основной результат, полученный в [2], можно сформулировать так: пусть  


имеется $n$-мерный параллелепипед $P_0$ с ребрами, параллельными 


координатным осям  


пространства $\bold R^n$ (занумерованным индексами $i=1,2,\dotsc,n$), 


составленный из  


параллелепипедов $P_1,P_2,\dotsc,P_m$, и на каждом из последних  


выбрано ровно по одному ребру, тогда  


найдется такой номер $i=1,2,\dotsc,n$, что ребро параллелепипеда $P_0$, 


параллельное $i$-oй 


координатной оси, выражается в виде линейной комбинации с коэффициентами 


$0,\,\pm1$ 


только тех из выбранных ребер параллелепипедов $P_1,P_2,\dotsc,P_m$, 


которые также  


параллельны $i$-ой координатной оси пространства $\bold R^n$. 


 


В данной работе эти соотношения между ребрами распространяются, с 


некоторыми  


оговорками, на грани параллелепипедов. Точнее, вместо ребер, 


соответствующих  


координатным осям пространства $\bold R^n$, берутся грани 


параллелепипедов,  


соответствующие разбиению множества индексов $\{1,2,\dotsc,n\}$ 


на произвольные его подмножества 


$\{N_1,N_2,\dotsc,N_k\}$.\footnote{Таким 


образом, приведенное  


выше утверждение для ребер соответствует значениям: $n=k$, 


$N_1=\{1\}$, $N_2=\{2\},\dotsc,N_n=\{n\}$.} 


В этом случае справедливо утверждение (теорема 2), аналогичное  


вышеприведенному, с той однако оговоркой, что коэффициенты в линейной  


комбинации могут быть произвольными рациональными числами, 


не обязательно равными $0,\,\pm1$. 


 


Доказательство этого утверждения является несложным следствием приводимой  


ниже (теорема 1) алгебраической характеристики  


пространства $n$-мерных параллелепипедов  


через пространства одномерных параллелепипедов. Точным формулировкам  


предпошлем ряд предварительных соглашений. 


 


Будем предполагать, что все рассматриваемые $n$-мерные параллелепипеды 


лежат в  


пространстве $\bold R^n_+=\{(t_1,t_2,\dotsc,t_n)\in\bold R^n:t_i\ge0,\ 


i=1,2,\dotsc,n\}$ и определяются системой неравенств вида 


\begin{equation} 


a_i\le t_i\le b_i,\quad i=1,2,\dotsc,n. 


\end{equation} 


Это техническое соглашение не ограничивает общности доказываемых ниже  


утверждений.  


 


Через $E^n$ обозначим векторное пространство над полем рациональных чисел 


$\bold Q$, 


порожденное характеристическими функциями\footnote{Под характеристической  


функцией параллелепипеда $P$, как обычно, понимаем функцию  


$\chi_P(t_1,t_2,\dotsc,t_n)=1$, если $(t_1,t_2,\dotsc,t_n)\in P$, и 


$\chi_P(t_1,t_2,\dotsc,t_n)=0$ в противном случае.} 


$\chi_P(t_1,t_2,\dotsc,t_n)$ параллелепипедов $P\subset\bold R^n_+$. 


Для краткости записи характеристические функции будем отождествлять с 


самими  


параллелепипедами и писать $P$ вместо $\chi_P(t_1,t_2,\dotsc,t_n)$. 


 


Удобно считать характеристические функции определенными с точностью  


до множества нулевого гиперобъема. Тогда тот факт, что  


параллелепипед $P_0$ составлен  


из параллелепипедов $P_1,P_2,\dotsc,P_m$, крaтко можно записать так: 


$P_0=P_1+P_2+\dotsb+P_m$, поскольку  


в этом случае равенство 


$\chi_{P_0}(t_1,t_2,\dotsc,t_n)=\chi_{P_1}(t_1,t_2,\dotsc,t_n)+ 


\chi_{P_2}(t_1,t_2,\dotsc,t_n)+\dotsb+\chi_{P_m}(t_1,t_2,\dotsc,t_n)$ 


имеет место всюду в $\bold R^n_+$, за исключением множества граничных 


точек  


параллелепипедов, т.\,е. за исключением множества нулевого гиперобъема.  


 


Для произвольной точки $(t_1,t_2,\dotsc,t_n)\in\bold R^n_+$ обозначим 


через $As(t_1,t_2,\dotsc,t_n)$ параллелепипед,  


определяемый системой неравенств $0\le u_i\le t_i$, $i=1,2,\dotsc,n$. 


Имеет место  


\medskip 


 


{\bf Лемма} [2]. {\it 


Семейство $\beta_n=\{As(t_1,t_2,\dotsc,t_n): t_i>0,\ i=1,2,\dotsc,n\}$ 


образует базис векторного  


пространства $E^n$, причем произвольный параллелепипед $P\in E^n$, 


определяемый системой  


$(1)$, следующим образом разлагается по элементам этого базиса\/$:$ 


\begin{equation} 


P=\sum(-1)^{p_1+p_2+\dotsb+p_n}As(c_1,c_2,\dotsc,c_n), 


\end{equation} 


где суммирование идет по всем $2^n$ вершинам  


параллелепипеда $P$, т.\,е. при каждом $i=1,2,\dotsc,n$ 


переменные $c_i$ независимо принимают значения $a_i$ или $b_i$,  


причем $p_i=1$ в случае, если $c_i=a_i$, и $p_i=0$ в случае, когда 


$c_i=b_i$. 


} 


\medskip 


 


{\bf Доказательство.} Покажем, что семейство $\beta_n$ линейно 


независимо в  


пространстве $E^n$. Пусть имеется $m$ различных элементов  


семейства $\beta_n$, удовлетворяющих равенству: $\lambda_1As(\mr{x}_1)+ 


\lambda_2As(\mr{x}_2)+\dotsb


+\lambda_m As(\mr{x}_m)=0$, где $\mr{x_1,x_2,\dotsc,x}_m$ --- 


различные точки пространства $\bold R^n_+$. 


Определим в пространстве $\bold R^n_+$ отношение порядка, положив 


$(t_1,t_2,\dotsc,t_n)\le(u_1,u_2,\dotsc,u_n)$ 


в случае, если $u_i\le t_i$ при всех $i=1,2,\dotsc,n$. Выберем в 


множестве $\{\mr{x_1,x_2,\dotsc,x}_m\}$ какой-нибудь  


максимальный элемент $\mr{x}_j$, тогда $\lambda_j=0$, т.\,к. разность множеств 


$As(\mr{x}_m)\setminus\cup_{k\ne j}As(\mr{x}_k)\ne\emptyset$, 


поскольку все точки $\mr{x_1,x_2,\dotsc,x}_m$ различны. Отсюда по индукции  


$\lambda_1=\lambda_2=\dotsb=\lambda_m=0$, что доказывает линейную 


независимость семейства $\beta_n$. 


 


Для доказательства равенства (2)  


возьмем произвольную точку $\mr{x}=(x_1,x_2,\dotsc,x_n)$ пространства $\bold 


R^n_+$. Возможны следующие  


три случая расположения точки $\mr{x}$ относительно параллелепипеда  


$P$. 


 


Точка $\mr{x}$ --- внутренняя точка параллелепипеда $P$. 


В этом случае $\mr{x}\in As(c_1,c_2,\dotsc,c_n)$ тогда  


и только тогда, когда $c_1=b_1,c_2=b_2,\dotsc,c_n=b_n$ и, следовательно, 


обе части равенства (2)  


принимают значение 1. 


 


B случае, когда $\mr{x}=(x_1,x_2,\dotsc,x_n)\notin P$, причем для 


некоторого индекса $i$ имеет место  


$x_i>b_i$, обе части равенства (2), очевидно, обращаются в нуль. 


 


Остается третий случай: $\mr{x}=(x_1,x_2,\dotsc,x_n)\notin P$, причем 


$x_i\le b_i$ при всех $i=1,2,\dotsc,n$ и $x_k<a_k$ 


для некоторого индекса $k$. В этом случае для произвольного набора 


$c_1,c_2,\dotsc,c_{k-1},c_{k+1},\dotsc,c_n$, 


где $c_i=a_i$ или $c_i=b_i$, $i=1,2,\dotsc,k-1,k+1,\dotsc,n$, либо 


$\mr{x}\notin As(c_1,c_2,\dotsc,c_{k-1},b_k,c_{k+1},\dotsc,c_n)$ и  


$\mr{x}\notin As(c_1,c_2,\dotsc,c_{k-1},a_k,c_{k+1},\dotsc,c_n)$, либо точка 


$\mr{x}$ принадлежит обоим этим параллелепипедам.  


Но названные параллелепипеды входят в правую часть равенства (2) с  


противоположными знаками, поэтому она обращается в нуль, 


как и левая часть (2). Таким образом, 


равенство (2), а вместе с ним и лемма, доказаны. 


 


Следующая теорема устанавливает, что пространство $E^n$ изоморфно 


тензорному  


произведению $n$ экземпляров пространства $E^1$. При этом изоморфизме  


параллелепипеду соответствует тензорное произведение его ребер.  


 


\begin{thm} 


Соответствие $As(t_1,t_2,\dotsc,t_n)\to As(t_1)\otimes As 


(t_2)\otimes\dotsb\otimes As(t_n)$, $t_i>0$, $i=1,2,\dotsc,n$, 


продолжается по линейности до изоморфизма векторного  


пространства $E^n$ в тензорное  


произведение $E^1\otimes E^1\otimes\dotsb\otimes E^1$\; $n$ векторных 


пространств, каждое из которых совпадает с 


$E^1$. При этом изоморфизме параллелепипеду, заданному системой $(1)$, 


соответствует произведение 


$(As(b_1)-As(a_1))\otimes (As(b_2)-As(a_2))\otimes \dotsb\otimes 


(As(b_n)-As(a_n))$. 


\end{thm} 


 


\begin{pf} 


Продолжив по линейности отображение $As(t_1,t_2,\dotsc,t_n)\to 


As(t_1)\otimes\linebreak As(t_2)\otimes \dotsb\otimes As(t_n)$, 


определенное на базисе $\beta_n$ векторного пространства $E^n$ (см.  


лемму), на все пространство $E^n$, получим  


линейное отображение $\varphi:E^n\to E^1\otimes E^1\otimes 


\dotsb\otimes E^1$. Нетрудно видеть, что так определенное  


отображение $\varphi$ является взаимно однозначным.  


Действительно, в силу доказанной леммы семейство 


$\beta_1=\{As(t):t>0\}$ является базисом  


пространства $E^1$. Следовательно, семейство 


$\{As(t_1)\otimes As(t_2)\otimes \dotsb\otimes As(t_n): t_i>0,\ 


i=1,2,\dotsc,n\}$ 


является базисом векторного пространства $E^1\otimes E^1\otimes 


\dotsb\otimes E^1$ [3] и, следовательно, $\varphi$ --- 


изоморфизм пространства $E^n$ на тензорное произведение $E^1\otimes 


E^1\otimes \dotsb\otimes E^1$. 


 


Пусть теперь $P$ --- параллелепипед, определяемый системой (1). Покажем, 


что  


$\varphi(P)=(As(b_1)-As(a_1))\otimes (As(b_2)-As(a_2))\otimes 


\dotsb\otimes (As(b_n)-As(a_n))$. Действительно, из равенства (2) имеем 


\begin{multline} 


\varphi(P)=\sum(-1)^{p_1+p_2+\dotsb+p_n}\varphi 


(As(c_1,c_2,\dotsc,c_n))= \\ =\sum 


(-1)^{p_1+p_2+\dotsb+p_n}As(c_1)\otimes 


As(c_2)\otimes\dotsb\otimes As(c_n), 


\end{multline} 


где $c_i$ независимо принимают значения $a_i$ или $b_i$ при каждом 


$i=1,2,\dotsc,n$, причем $p_i=1$ в  


случае, если $c_i=a_i$, и $p_i=0$ в случае, когда $c_i=b_i$. С другой 


стороны, раскрывая скобки в  


произведении $(As(b_1)-As(a_1))\otimes (As(b_2)-As(a_2))\otimes 


\dotsb\otimes (As(b_n)-As(a_n))$, нетрудно видеть, что  


оно совпадает с правой частью равенства (3). 


\end{pf} 


 


Tеорема 1 позволяет распространить приведенное в начале статьи  


утверждение о соотношениях между ребрами параллелепипедов на  


их грани следующим образом. 


 


\begin{thm} 


Пусть $n$-мерный параллелепипед $P_0\subset\bold R^n$ составлен из 


параллелепипедов 


$P_1,P_2,\dotsc,P_m$ и зафиксировано какое-либо разбиение множества 


координатных осей 


$N=\{1,2,\dotsc,n\}$ пространства $\bold R^n$ на $k$ подмножеств\/$:$ 


$N_1,N_2,\dotsc,N_k$. Тогда для любого покрытия  


множества индексов $M=\{1,2,\dotsc,m\}$ его произвольными 


подмножествами $M_1,M_2,\dotsc,M_k$ 


найдется номер $i$ такой, что при некоторых рациональных $\varepsilon_j$ 


имеет место  


равенство 


\begin{equation} 


P^i_0=\sum_{j\in M_i}\varepsilon_j P^i_j, 


\end{equation} 


где $P^i_j$, $j=0,1,\dotsc,m$, --- проекция параллелепипеда $P_j$ на 


произведение $\bold R^{N_i}$ тех  


координатных осей пространства $\bold R^n$, 


номера которых входят в множество индексов $N_i$.  


\end{thm} 


 


{\bf Доказательство} проведем от противного. Пусть имеется какое-нибудь  


разбиение множества координатных осей $N=\{1,2,\dotsc,n\}$ на $k$ 


подмножеств 


$N_1,N_2,\dotsc,N_k$ и произвольное покрытие множества индексов 


$M=\{1,2,\dotsc,m\}$ его  


подмножествами $M_1,M_2,\dotsc,M_k$. 


Предположим, что для любого $i=1,2,\dotsc,k$ равенство (4) не имеет 


места ни при каких  


рациональных значениях $\varepsilon_j$. Это означает [3], что для 


каждого $i=1,2,\dotsc,k$ найдется  


линейная форма $F_i:E^{N_i}\to\bold Q$ такая, что 


\begin{equation} 


F_i(P^i_0)=1,\ \; 


\text{а}\ \; F_i(P^i_j)=0 


\end{equation} 


для всех $j\in M_i$, где $E^{N_i}$ --- векторное пространство,  


порожденное параллелепипедами, лежащими в пространстве $\bold R^{N_i}$. 


 


В силу свойств тензорного произведения [3] отображение $F$ пространства 


$E^n$: 


\begin{equation} 


F(As(t_1)\otimes As(t_2)\otimes\dotsb\otimes As(t_n))= 


\prod^k_{i=1}F_i(As(t_{i_1})\otimes As(t_{i_2}) 


\otimes\dotsb\otimes As(t_{i_r})), 


\end{equation}  


где $(t_1,t_2,\dotsc,t_n)\in\bold R^n_+$, a $\{i_1,i_2,\dotsc,i_r\}\in 


N_i$, является линейной формой на пространстве $E^n$. 


 


Найдем значения $F(P_j)$, $j=0,1,2,\dotsc,m$. В силу (6) и теоремы 1 


$F(P_j)=\prod\limits^k_{i=1}F_i(P^i_j)$. Отсюда $F(P_0)=1$, поскольку 


$F_i(P^i_0)=1$, $i=1,2,\dotsc,k$, в силу равенства (5). 


При $j=1,2,\dotsc,m$ из того же равенства (5) имеем $F(P_j)=0$, т.\,к. 


$F_i(P^i_j)=0$ для того  


индекса $i$, при котором $j\in M_i$. 


 


Найденные значения $F(P_j)$, $j=0,1,2,\dotsc,m$, противоречат условию 


теоремы, поскольку  


если параллелепипед $P_0$ составлен из параллелепипедов 


$P_1,P_2,\dotsc,P_m$, т.\,е.  


$P_0=P_1+P_2+\dotsb+P_m$, то для линейной формы $F$ должно выполняться 


равенство $F(P_0)=F(P_1)+F(P_2)+\dotsb\linebreak +F(P_m)$. $\quad\square$


 


В заключение по поводу вышеизложенного сделаем два замечания.


 


\begin{note} 


Утверждение теоремы 2 сохраняет силу и при следующем более слабом, нежели 


приведенное в ее формулировке, предположении:  


параллелепипед $P_0$ является линейной  


комбинацией с произвольными коэффициентами  


параллелепипедов $P_1,P_2,\dotsc,P_m$, т.\,е.  


$P_0=\lambda_1P_1+\lambda_2P_2+\dotsb+\lambda_m P_m$. Нетрудно видеть, 


что приведенное доказательство сохранит силу  


и в этом случае.  


\end{note} 


 


\begin{note} 


Kак показывают довольно громоздкие примеры, для некоторых разбиений  


равенство (4) может иметь место лишь при дробных  


значениях коэффициентов $\varepsilon_j$ в отличие от соответствующего  


утверждения для ребер, приведенного в начале статьи (см. также  


[2]). 


\end{note} 


 


\begin{thebibliography}{9} 


 


\bibitem{1} 


Вarnes F. {\em Algebraic theory of brick packing}. II // 


Discrete Math. -- 1982. -- V.\,42. -- \No\,2--3. -- P.\,129--144. 


\bibitem{2} 


Покровский В.Г. {\em О линейных соотношениях при разрезаниях на  


$n$-мерные параллелепипеды} // Матем. заметки. -- 1985. -- T.\,37. 


-- \No\,6. -- C.\,901--907. 


\bibitem{3} 


Ленг С. {\em Алгебра}. -- М.: Мир, 1968. -- 564 с. 


 


\end{thebibliography} 


 


\vskip 2cm 


 


\post{\it Московский городской}{\it Поступила} 


\post{\it педагогический университет}{20.06.2000} 


 


\label{end} 


\end{document} 


